% Appendix fragment for `main (6).tex`.
% Required preamble additions:
%   \usepackage{longtable}
%   \usetikzlibrary{arrows.meta,positioning,fit}

\section{Machine-checked formalization of Sections 4 and 5}
\label{app:lean-formalization}

\subsection{Scope and trust boundary}
\label{app:lean-scope}

The source currently used for Sections~4 and~5 is the attached formalization
blueprint: the corresponding headings in the main manuscript contain only the
sentence ``See attached note.''  Consequently, the correspondence table below
uses the stable labels and titles in that blueprint.  Its sufficiency and
necessity parts contain 152 labelled formal environments (55 definitions,
73 lemmas, 18 propositions, four corollaries, and two theorems).

The accompanying development is checked with Lean~4.32.2 and
mathlib~4.32.2.  It contains no \texttt{sorry}, \texttt{admit},
project-defined \texttt{axiom}, or project-defined \texttt{opaque}
declaration.  The command
\begin{verbatim}
powershell -NoProfile -ExecutionPolicy Bypass -File scripts\check.ps1
\end{verbatim}
checks both the source and the kernel dependencies of the exported results.
The closed declaration \texttt{verifiedKernelBundle} collects the principal
finite algebra, order, stationarity, midpoint, and dominant-block results.

This is a coverage report, not a claim that the classification theorem is
already machine checked.  In particular, the present Lean source does not
declare \texttt{mainClassification}, \texttt{necessityBundle}, or
\texttt{allConditionalEntropyAxioms}.  Those declarations would be
misleading until the endpoint-aware R\'enyi library, measure
discretization, differentiated signed integrals, and the full block and
Shannon localization arguments have been translated.  The table records
these obligations explicitly instead of representing them by hypotheses or
project-specific axioms.

Two typing conventions in the blueprint should also be kept visible.  First,
the singular moment at the Shannon point is represented by one-sided
extended-real moments plus a finiteness guard; an unguarded ordinary signed
integral need not exist.  Second, the Lean-facing blueprint uses natural
logarithms, so the fair-bit normalization is \(\log 2\); the normalization
equal to one in the main manuscript is obtained by division by this positive
constant.

\subsection{Dependency graph}
\label{app:lean-dependency-graph}

Figure~\ref{fig:lean-dependency-graph} is the readable module-layer dependency
graph.  A blue box contains at least one
kernel-checked component but is not necessarily complete.  An orange box is
specified by the blueprint but still lacks a faithful Lean translation.

\begin{figure}[htbp]
  \centering
  \begin{tikzpicture}[
      x=1cm,y=1cm,
      dep/.style={-{Stealth[length=1.7mm]},draw=black!55,line width=.45pt},
      partial/.style={draw=blue!65!black,fill=blue!8,rounded corners=2pt,
        align=center,font=\scriptsize,minimum height=7mm,text width=38mm},
      open/.style={draw=orange!80!black,fill=orange!10,rounded corners=2pt,
        align=center,font=\scriptsize,minimum height=7mm,text width=38mm}]
    \node[partial] (m1) at (-3.1,0) {M1. Finite carriers, cones, and channels\\\emph{partial}};
    \node[open]    (m2) at ( 3.1,0) {M2. Compact parameters, measures, support, and moments\\\emph{open}};
    \node[open]    (m3) at (0,-1.25) {M3. Total endpoint-aware R\'enyi library\\\emph{open}};
    \node[partial] (m4) at (0,-2.50) {M4. Conditional candidates and shape reductions\\\emph{partial}};
    \node[partial] (m5) at (-3.1,-3.75) {M5. Temperate sufficiency\\\emph{partial algebraic kernels}};
    \node[open]    (m6) at (-3.1,-5.00) {M6. Tropical and derivation sufficiency\\\emph{open}};
    \node[partial] (m7) at ( 3.1,-3.75) {M7. Differentiation and signed integration\\\emph{partial stationarity kernels}};
    \node[partial] (m8) at ( 3.1,-5.00) {M8. Two- and three-block localization\\\emph{partial dominant-block bounds}};
    \node[open]    (m9) at ( 3.1,-6.25) {M9. Shannon-point localization\\\emph{open}};
    \node[partial] (m10) at (0,-7.50) {M10. Necessity branches\\\emph{partial scalar kernels}};
    \node[open]    (m11) at (0,-8.75) {M11. Classification and embedded-axiom lift\\\emph{open}};

    \draw[dep] (m1) -- (m3);
    \draw[dep] (m2) -- (m3);
    \draw[dep] (m1) -- (m4);
    \draw[dep] (m2) -- (m4);
    \draw[dep] (m3) -- (m4);
    \draw[dep] (m2) -- (m5);
    \draw[dep] (m3) -- (m5);
    \draw[dep] (m4) -- (m5);
    \draw[dep] (m3) -- (m6);
    \draw[dep] (m4) -- (m6);
    \draw[dep] (m5) -- (m6);
    \draw[dep] (m2) -- (m7);
    \draw[dep] (m3) -- (m7);
    \draw[dep] (m4) -- (m7);
    \draw[dep] (m1) -- (m8);
    \draw[dep] (m2) -- (m8);
    \draw[dep] (m3) -- (m8);
    \draw[dep] (m7) -- (m8);
    \draw[dep] (m7) -- (m9);
    \draw[dep] (m8) -- (m9);
    \draw[dep] (m4) -- (m10);
    \draw[dep] (m7) -- (m10);
    \draw[dep] (m8) -- (m10);
    \draw[dep] (m9) -- (m10);
    \draw[dep] (m4) -- (m11);
    \draw[dep] (m5) -- (m11);
    \draw[dep] (m6) -- (m11);
    \draw[dep] (m10) -- (m11);
  \end{tikzpicture}
  \caption{Condensed formalization dependency graph.  An arrow
  \(A\to B\) means that blueprint layer \(B\) is planned to depend on
  layer \(A\).  Colours report translation status, not mathematical truth.}
  \label{fig:lean-dependency-graph}
\end{figure}

The repository also contains Graphviz and SVG versions of the same condensed
graph.  A separate declaration-level dependency index retains the individual
Lean names that are intentionally omitted from
Figure~\ref{fig:lean-dependency-graph}.

\subsection{Statement correspondence}
\label{app:lean-correspondence}

The status ``proved and kernel checked'' means that the named Lean declaration
has the blueprint statement's full conclusion.  ``Partial algebraic kernel''
means that it checks a genuine calculation used by the manuscript statement
but does not prove that statement's full measure-theoretic or asymptotic
conclusion.  ``Not formalized'' means that no Lean theorem with the manuscript's
conclusion is present.  Definitions are marked separately from proved
propositions.

% Statement ledger for the blueprint's Sections 4--5 and final assembly.
% This fragment is included by lean-formalization-appendix.tex.

\providecommand{\BlueprintStatement}[2]{%
  \texttt{\detokenize{#1}}: \emph{#2}}
\providecommand{\LeanDecl}[1]{\texttt{\detokenize{#1}}}

\small
\begin{longtable}{>{\raggedright\arraybackslash}p{0.35\textwidth}
                  >{\raggedright\arraybackslash}p{0.42\textwidth}
                  >{\raggedright\arraybackslash}p{0.15\textwidth}}
\caption{Statement-by-statement correspondence for blueprint Sections 4--5
and final assembly.}
\label{tab:blueprint-statement-correspondence}\\
\toprule
Manuscript statement & Lean declaration & Status\\
\midrule
\endfirsthead
\toprule
Manuscript statement & Lean declaration & Status\\
\midrule
\endhead
\bottomrule
\endfoot

\BlueprintStatement{lem:origin-extension}{Extension from the punctured cone}
  & \LeanDecl{ConditionalEntropy.originExtension} & Proved and kernel checked\\
\BlueprintStatement{lem:punctured-cone-convex}{Convexity of the punctured nonnegative cone}
  & \LeanDecl{ConditionalEntropy.puncturedConeConvex} & Proved and kernel checked\\
\BlueprintStatement{lem:composition-rule}{Monotone composition rule}
  & \LeanDecl{ConditionalEntropy.compositionRule} & Proved and kernel checked\\
\BlueprintStatement{lem:cone-composition-rule}{Monotone composition on the punctured cone}
  & \LeanDecl{ConditionalEntropy.coneCompositionRule} & Proved and kernel checked\\
\BlueprintStatement{lem:finite-holder}{Finite H\"older inequality}
  & --- & Not formalized\\
\BlueprintStatement{lem:power-curvature}{Power-mean curvature}
  & --- & Not formalized\\
\BlueprintStatement{lem:parameter-power-curvature}{Compactified power-mean bridge}
  & --- & Not formalized\\
\BlueprintStatement{def:hessian-quad}{Hessian quadratic form}
  & --- & Not formalized\\
\BlueprintStatement{lem:hessian-sign-criterion}{Hessian sign criterion}
  & --- & Not formalized\\
\BlueprintStatement{lem:real-rpow-calculus}{Positive-base real-power calculus}
  & --- & Not formalized\\
\BlueprintStatement{lem:monomial-curvature}{Curvature of an affine monomial}
  & \LeanDecl{ConditionalEntropy.monomialCurvature};
    \LeanDecl{ConditionalEntropy.monomialCurvature_eq_neg_variance};
    \LeanDecl{ConditionalEntropy.monomialCurvature_nonpos_of_nonneg};
    \LeanDecl{ConditionalEntropy.coefficient_mem_Icc_of_curvature_nonpos};
    \LeanDecl{ConditionalEntropy.monomialCurvature_pos_of_negative_coefficient};
    \LeanDecl{ConditionalEntropy.two_positive_stationary_witness};
    \LeanDecl{ConditionalEntropy.negativeTemperate_two_coefficient_curvature}
  & Partial algebraic kernel\\
\BlueprintStatement{lem:finite-support-enumeration}{Enumeration of a finite probability support}
  & --- & Not formalized\\
\BlueprintStatement{lem:finite-support-dirac}{Finite-support measure decomposition}
  & --- & Not formalized\\
\BlueprintStatement{prop:finite-sufficiency}{Finite-support sufficiency}
  & \LeanDecl{ConditionalEntropy.positiveTemperate_coefficients_nonnegative};
    \LeanDecl{ConditionalEntropy.negativeTemperate_two_coefficient_curvature}
  & Partial algebraic kernel\\
\BlueprintStatement{prop:general-sufficiency}{General temperate sufficiency}
  & \LeanDecl{ConditionalEntropy.isConcave_of_pointwise_tendsto};
    \LeanDecl{ConditionalEntropy.isConvex_of_pointwise_tendsto};
    \LeanDecl{ConditionalEntropy.conditioning_monotone_of_concave};
    \LeanDecl{ConditionalEntropy.conditioning_antitone_of_convex}
  & Partial algebraic kernel\\
\BlueprintStatement{cor:finite-t-sufficiency}{Finite-$t$ sufficiency}
  & \LeanDecl{ConditionalEntropy.positiveTemperate_coefficients_nonnegative};
    \LeanDecl{ConditionalEntropy.negativeTemperate_two_coefficient_curvature};
    \LeanDecl{ConditionalEntropy.conditioning_monotone_of_concave};
    \LeanDecl{ConditionalEntropy.conditioning_antitone_of_convex}
  & Partial algebraic kernel\\
\BlueprintStatement{prop:negative-tropical-sufficiency}{Negative tropical sufficiency}
  & \LeanDecl{ConditionalEntropy.isQuasiConvex_of_pointwise_tendsto}
  & Partial algebraic kernel\\
\BlueprintStatement{lem:positive-tropical-component}{Positive tropical component formula}
  & --- & Not formalized\\
\BlueprintStatement{prop:positive-tropical-sufficiency}{Positive tropical sufficiency}
  & \LeanDecl{ConditionalEntropy.isStronglyQuasiConcave_of_pointwise_tendsto}
  & Partial algebraic kernel\\
\BlueprintStatement{def:derivation-family}{Derivation polymorphic family}
  & --- & Not formalized\\
\BlueprintStatement{prop:derivation-sufficiency}{Derivation sufficiency}
  & --- & Not formalized\\
\BlueprintStatement{def:signed-column-functions}{Signed-measure column functions}
  & --- & Not formalized\\
\BlueprintStatement{def:signed-witness-predicates}{Positive and negative signed witnesses}
  & --- & Not formalized\\
\BlueprintStatement{lem:signed-witness-bridge}{Scaling, total variation, support, and atoms of witnesses}
  & --- & Not formalized\\
\BlueprintStatement{lem:signed-scalar-column-bridge}{Signed-scalar bridge to the candidate column functions}
  & --- & Not formalized\\
\BlueprintStatement{def:finite-param}{Embedding a finite real order}
  & --- & Not formalized\\
\BlueprintStatement{def:positive-line-data}{Positive multiplicative line data}
  & --- & Not formalized\\
\BlueprintStatement{def:positive-line-raw}{Raw multiplicative line}
  & --- & Not formalized\\
\BlueprintStatement{def:positive-line-predicate}{Positivity predicate for a line}
  & --- & Not formalized\\
\BlueprintStatement{lem:line-positive-zero}{The base point of a positive line is positive}
  & --- & Not formalized\\
\BlueprintStatement{def:line-cone}{Proof-carrying cone versions of a positive line}
  & --- & Not formalized\\
\BlueprintStatement{def:positive-line-prob}{Total normalized line}
  & --- & Not formalized\\
\BlueprintStatement{def:scalar-qcvx-on}{Scalar quasiconvexity on a set}
  & \LeanDecl{ConditionalEntropy.IsQuasiConvex}
  & Partial algebraic kernel\\
\BlueprintStatement{lem:cone-affine-line-bridge}{Restricting cone curvature to a positive affine line}
  & --- & Not formalized\\
\BlueprintStatement{def:entropy-line}{Entropy line}
  & --- & Not formalized\\
\BlueprintStatement{def:entropy-line-first}{First entropy line derivative}
  & --- & Not formalized\\
\BlueprintStatement{def:entropy-line-second}{Second entropy line derivative}
  & --- & Not formalized\\
\BlueprintStatement{def:escort-weight}{Escort weight}
  & --- & Not formalized\\
\BlueprintStatement{def:effective-line-velocity}{Effective line velocity}
  & --- & Not formalized\\
\BlueprintStatement{def:escort-mean}{Escort mean}
  & \LeanDecl{ConditionalEntropy.weightedMean}
  & Partial algebraic kernel\\
\BlueprintStatement{def:escort-second}{Escort second moment}
  & \LeanDecl{ConditionalEntropy.weightedSecondMoment}
  & Partial algebraic kernel\\
\BlueprintStatement{def:escort-variance}{Escort variance}
  & \LeanDecl{ConditionalEntropy.weightedVariance}
  & Partial algebraic kernel\\
\BlueprintStatement{def:fixed-max-coordinate}{Fixed maximal-coordinate condition}
  & --- & Not formalized\\
\BlueprintStatement{lem:positivity-intervals}{Common positivity and maximal-coordinate intervals}
  & --- & Not formalized\\
\BlueprintStatement{def:iterated-deriv}{Iterated one-variable derivative}
  & --- & Not formalized\\
\BlueprintStatement{def:lambda-deriv}{Iterated derivative in the second variable}
  & --- & Not formalized\\
\BlueprintStatement{def:alpha-lambda-deriv}{Mixed parameter--line derivative}
  & --- & Not formalized\\
\BlueprintStatement{def:removable-shannon-quotient}{Removable Shannon quotient}
  & --- & Not formalized\\
\BlueprintStatement{lem:parameterized-removable-quotient}{Parameterized removable quotient at the Shannon point}
  & --- & Not formalized\\
\BlueprintStatement{lem:signed-integral-differentiate-twice}{Twice differentiating a finite signed-measure integral}
  & --- & Not formalized\\
\BlueprintStatement{def:shannon-line-slope}{Shannon line slope}
  & --- & Not formalized\\
\BlueprintStatement{lem:exact-derivatives}{Exact entropy derivatives}
  & --- & Not formalized\\
\BlueprintStatement{def:integrated-entropy-line}{Integrated entropy along a line}
  & --- & Not formalized\\
\BlueprintStatement{def:signed-log-phi-line}{Signed logarithmic column line}
  & --- & Not formalized\\
\BlueprintStatement{lem:signed-line-column-bridge}{Exact bridge from line kernels to signed column functions}
  & --- & Not formalized\\
\BlueprintStatement{lem:differentiate-integral}{Differentiation through the parameter integral}
  & --- & Not formalized\\
\BlueprintStatement{lem:null-thresholds}{Null thresholds}
  & --- & Not formalized\\
\BlueprintStatement{lem:two-upper-null-thresholds}{Two finite null thresholds between upper parameters}
  & --- & Not formalized\\
\BlueprintStatement{lem:support-strict-integral}{Support gives strict weighted mass}
  & --- & Not formalized\\
\BlueprintStatement{lem:scalar-measure-support}{Scalar measure support and null points}
  & --- & Not formalized\\
\BlueprintStatement{def:block-index}{Finite dependent block carrier}
  & --- & Not formalized\\
\BlueprintStatement{def:block-vector}{Constant-on-block vector}
  & --- & Not formalized\\
\BlueprintStatement{def:block-data}{Asymptotic block data}
  & --- & Not formalized\\
\BlueprintStatement{lem:block-count-pos}{Positive block counts}
  & --- & Not formalized\\
\BlueprintStatement{def:block-carrier-nonempty}{Nonempty block-carrier instance}
  & --- & Not formalized\\
\BlueprintStatement{def:block-index-nonempty}{Nonempty arbitrary block index}
  & --- & Not formalized\\
\BlueprintStatement{def:block-max}{Maximum of a nonempty block vector}
  & --- & Not formalized\\
\BlueprintStatement{def:block-line-data}{Typed block line}
  & --- & Not formalized\\
\BlueprintStatement{def:block-kernels}{Block exponents, escort weights, and derivative kernels}
  & --- & Not formalized\\
\BlueprintStatement{lem:block-vector-formulas}{Block-vector finite-sum formulas}
  & --- & Not formalized\\
\BlueprintStatement{lem:rounding}{Rounding estimates}
  & --- & Not formalized\\
\BlueprintStatement{lem:dominant-block}{Dominant-block escort estimate}
  & \LeanDecl{ConditionalEntropy.dominantBlock_first};
    \LeanDecl{ConditionalEntropy.dominantBlock_second}
  & Partial algebraic kernel\\
\BlueprintStatement{lem:near-shannon-dominant}{Near-Shannon cancellation for a dominant block}
  & --- & Not formalized\\
\BlueprintStatement{lem:large-alpha-tail}{Large-order tail bound}
  & --- & Not formalized\\
\BlueprintStatement{def:block-limit-kernels}{Pointwise limit kernels for a block family}
  & --- & Not formalized\\
\BlueprintStatement{prop:block-limit-passage}{Uniform passage to the block limit}
  & --- & Not formalized\\
\BlueprintStatement{def:special-block-data}{Specialized block data and signed truncations}
  & --- & Not formalized\\
\BlueprintStatement{def:special-block-top-unique}{Unique top blocks for the special families}
  & --- & Not formalized\\
\BlueprintStatement{def:block-dominance-maps}{Total two-block and three-block dominance maps}
  & --- & Not formalized\\
\BlueprintStatement{lem:two-block-dominance-package}{Exact two-block dominance package}
  & --- & Not formalized\\
\BlueprintStatement{lem:three-block-dominance-package}{Exact three-block dominance package}
  & --- & Not formalized\\
\BlueprintStatement{lem:two-order-one-dominance}{Unique order-one block for the two-block family}
  & --- & Not formalized\\
\BlueprintStatement{lem:three-order-one-dominance}{Unique order-one block for the three-block family}
  & --- & Not formalized\\
\BlueprintStatement{def:block-log-g-kernels}{Block logarithmic and norm-free kernels}
  & --- & Not formalized\\
\BlueprintStatement{lem:compact-uniform-eval}{Evaluation of compact-uniform convergence}
  & --- & Not formalized\\
\BlueprintStatement{lem:compact-uniform-add}{Addition closure for compact-uniform convergence}
  & --- & Not formalized\\
\BlueprintStatement{lem:compact-uniform-singleton}{Evaluation on a singleton}
  & --- & Not formalized\\
\BlueprintStatement{lem:subtype-uniform-error-bridge}{Subtype bridge for uniform signed-integral errors}
  & --- & Not formalized\\
\BlueprintStatement{lem:block-curve-kernel-bridge}{Block curve--kernel bridge}
  & --- & Not formalized\\
\BlueprintStatement{lem:block-g-kernel-integrals}{Block (G)-kernels are signed kernel integrals}
  & --- & Not formalized\\
\BlueprintStatement{lem:block-log-mass-limit}{Block logarithmic-mass splitting and limit}
  & --- & Not formalized\\
\BlueprintStatement{lem:block-log-mass-continuity}{Continuity of logarithmic block-mass kernels}
  & --- & Not formalized\\
\BlueprintStatement{lem:block-kernel-regularity}{Linearity and continuity of the block kernels}
  & --- & Not formalized\\
\BlueprintStatement{def:two-block-limits}{Two-block limit polynomials}
  & --- & Not formalized\\
\BlueprintStatement{def:three-block-limits}{Three-block cells and limit polynomials}
  & --- & Not formalized\\
\BlueprintStatement{lem:two-block-limit-integral-eval}{Evaluation of the two-block limit integrals}
  & --- & Not formalized\\
\BlueprintStatement{lem:three-block-limit-integral-eval}{Evaluation of the three-block limit integrals}
  & --- & Not formalized\\
\BlueprintStatement{prop:block-localisation}{Two-block and three-block localisation}
  & --- & Not formalized\\
\BlueprintStatement{prop:two-block-localisation}{Two-block localisation interface}
  & --- & Not formalized\\
\BlueprintStatement{prop:three-block-localisation}{Three-block localisation interface}
  & --- & Not formalized\\
\BlueprintStatement{def:shannon-data}{Shannon block data}
  & --- & Not formalized\\
\BlueprintStatement{lem:shannon-q-pos}{Positivity of the complementary Shannon weight}
  & --- & Not formalized\\
\BlueprintStatement{lem:shannon-log-scale-pos}{Positivity of the Shannon logarithmic scale}
  & --- & Not formalized\\
\BlueprintStatement{lem:shannon-count-pos}{Positive Shannon block counts}
  & --- & Not formalized\\
\BlueprintStatement{def:shannon-index-nonempty}{Nonempty Shannon index and block representative}
  & --- & Not formalized\\
\BlueprintStatement{def:shannon-line-data}{Shannon line and its escort}
  & --- & Not formalized\\
\BlueprintStatement{lem:shannon-top-unique-at-zero}{The third Shannon block is strictly maximal at the base point}
  & --- & Not formalized\\
\BlueprintStatement{def:shannon-kernels}{Named Shannon kernels}
  & --- & Not formalized\\
\BlueprintStatement{lem:shannon-curve-kernel-bridge}{Shannon curve--kernel bridge}
  & --- & Not formalized\\
\BlueprintStatement{lem:shannon-kernel-integral-bridge}{Shannon derivative--integral and logarithmic splitting bridge}
  & --- & Not formalized\\
\BlueprintStatement{lem:shannon-kernel-regularity}{Linearity and continuity of the Shannon kernels}
  & --- & Not formalized\\
\BlueprintStatement{lem:shannon-dedicated-escort}{Dedicated Shannon escort, compact-region, and tail estimates}
  & --- & Not formalized\\
\BlueprintStatement{def:shannon-remainder}{Shannon main term and uniform remainder}
  & --- & Not formalized\\
\BlueprintStatement{lem:uniform-shannon-expansion}{Uniform Shannon expansion}
  & --- & Not formalized\\
\BlueprintStatement{lem:shannon-neighbourhood}{Uniform control near the Shannon point}
  & --- & Not formalized\\
\BlueprintStatement{def:remove-signed-atom}{Signed atom and its removal}
  & --- & Not formalized\\
\BlueprintStatement{lem:remove-signed-atom}{Removing a signed atom}
  & --- & Not formalized\\
\BlueprintStatement{lem:signed-witness-integral-bridge}{Witness atoms and the positive lower truncation}
  & --- & Not formalized\\
\BlueprintStatement{lem:shannon-log-mass}{Shannon-block logarithmic mass derivatives}
  & --- & Not formalized\\
\BlueprintStatement{def:shannon-localization-targets}{Shannon localization targets}
  & --- & Not formalized\\
\BlueprintStatement{prop:shannon-localisation}{Shannon-point localisation}
  & --- & Not formalized\\
\BlueprintStatement{def:normalized-curvature}{Normalized scalar-line curvature}
  & --- & Not formalized\\
\BlueprintStatement{lem:log-curvature-identity}{Logarithmic curvature identity}
  & --- & Not formalized\\
\BlueprintStatement{lem:log-curvature-obstruction}{Logarithmic curvature limit obstruction}
  & --- & Not formalized\\
\BlueprintStatement{def:stationarity-package}{Stationarity convergence package}
  & --- & Not formalized\\
\BlueprintStatement{lem:stationarity-correction}{Exact stationarity correction}
  & \LeanDecl{ConditionalEntropy.stationarityCorrection_dot};
    \LeanDecl{ConditionalEntropy.stationarityCorrection_tendsto}
  & Partial algebraic kernel\\
\BlueprintStatement{lem:quasiconvex-second}{Stationary curvature of a quasi-convex function}
  & \LeanDecl{ConditionalEntropy.not_quasiConvex_of_strict_midpoint_peak}
  & Partial algebraic kernel\\
\BlueprintStatement{lem:typed-log-curvature-contradiction}{Typed curvature contradiction from logarithmic kernels}
  & \LeanDecl{ConditionalEntropy.not_concave_of_strict_midpoint_valley};
    \LeanDecl{ConditionalEntropy.not_convex_of_strict_midpoint_peak}
  & Partial algebraic kernel\\
\BlueprintStatement{lem:corrected-stationary-line-obstruction}{Corrected stationary-line obstruction}
  & \LeanDecl{ConditionalEntropy.stationarityCorrection_dot};
    \LeanDecl{ConditionalEntropy.stationarityCorrection_tendsto};
    \LeanDecl{ConditionalEntropy.not_quasiConvex_of_strict_midpoint_peak}
  & Partial algebraic kernel\\
\BlueprintStatement{prop:positive-necessity}{Positive-measure obstruction}
  & \LeanDecl{ConditionalEntropy.negative_tail_concavity_obstruction};
    \LeanDecl{ConditionalEntropy.excessive_lower_moment_concavity_obstruction};
    \LeanDecl{ConditionalEntropy.shannon_atom_concavity_obstruction};
    \LeanDecl{ConditionalEntropy.positiveTemperate_shannon_atom_zero};
    \LeanDecl{ConditionalEntropy.positiveTemperate_no_upper_moment};
    \LeanDecl{ConditionalEntropy.positiveTemperate_lowerMoment_le_one};
    \LeanDecl{ConditionalEntropy.positiveTemperate_necessary}
  & Partial algebraic kernel\\
\BlueprintStatement{cor:positive-temperate-probability-necessity}{Positive temperate necessity for probability measures}
  & \LeanDecl{ConditionalEntropy.positiveTemperate_necessary}
  & Partial algebraic kernel\\
\BlueprintStatement{prop:negative-shannon-obstruction}{A Shannon atom excludes upper support}
  & --- & Not formalized\\
\BlueprintStatement{def:negative-truncations}{Positive lower and upper truncations}
  & --- & Not formalized\\
\BlueprintStatement{lem:negative-truncation-bridge}{Negative-witness truncation identities}
  & --- & Not formalized\\
\BlueprintStatement{prop:truncated-moment}{Truncated-moment obstruction}
  & \LeanDecl{ConditionalEntropy.two_positive_stationary_witness};
    \LeanDecl{ConditionalEntropy.negativeTemperate_truncated_moment};
    \LeanDecl{ConditionalEntropy.negativeTemperate_exceptional_moment}
  & Partial algebraic kernel\\
\BlueprintStatement{prop:one-upper-point}{At most one upper support point}
  & \LeanDecl{ConditionalEntropy.two_positive_stationary_witness};
    \LeanDecl{ConditionalEntropy.negativeTemperate_atMostOneUpperCell}
  & Partial algebraic kernel\\
\BlueprintStatement{prop:negative-temperate-necessity}{Complete negative temperate necessity}
  & \LeanDecl{ConditionalEntropy.negativeTemperate_exceptional_moment};
    \LeanDecl{ConditionalEntropy.negativeTemperate_atMostOneUpperCell}
  & Partial algebraic kernel\\
\BlueprintStatement{cor:negative-temperate-probability-necessity}{Negative temperate necessity for probability measures}
  & \LeanDecl{ConditionalEntropy.negativeTemperate_exceptional_moment};
    \LeanDecl{ConditionalEntropy.negativeTemperate_atMostOneUpperCell}
  & Partial algebraic kernel\\
\BlueprintStatement{prop:negative-tropical-necessity}{Complete negative tropical necessity}
  & \LeanDecl{ConditionalEntropy.negativeTropical_moment_nonnegative};
    \LeanDecl{ConditionalEntropy.negativeTropical_exceptional_moment}
  & Partial algebraic kernel\\
\BlueprintStatement{def:prob-as-pos-cone}{Probability support as a positive cone vector}
  & --- & Not formalized\\
\BlueprintStatement{lem:prob-as-pos-cone-identities}{Probability-to-cone identities}
  & --- & Not formalized\\
\BlueprintStatement{lem:prob-support-nonempty}{A probability vector has nonempty positive support}
  & --- & Not formalized\\
\BlueprintStatement{lem:same-support-decomposition}{Decomposition inside a simplex face}
  & --- & Not formalized\\
\BlueprintStatement{lem:probability-eq-dirac-zero}{Probability supported at zero is the zero Dirac measure}
  & --- & Not formalized\\
\BlueprintStatement{prop:positive-tropical-necessity}{Only the support-rank measure survives}
  & --- & Not formalized\\
\BlueprintStatement{def:derivation-column-function}{Derivation column function}
  & --- & Not formalized\\
\BlueprintStatement{prop:derivation-necessity}{Derivation necessity}
  & \LeanDecl{ConditionalEntropy.positive_tail_derivation_obstruction};
    \LeanDecl{ConditionalEntropy.derivation_no_positive_upper_tail};
    \LeanDecl{ConditionalEntropy.derivation_necessary}
  & Partial algebraic kernel\\
\BlueprintStatement{thm:necessity-bundle}{All necessary parameter conditions}
  & \LeanDecl{ConditionalEntropy.positiveTemperate_necessary};
    \LeanDecl{ConditionalEntropy.negativeTemperate_exceptional_moment};
    \LeanDecl{ConditionalEntropy.negativeTropical_exceptional_moment};
    \LeanDecl{ConditionalEntropy.derivation_necessary}
  & Partial algebraic kernel\\
\BlueprintStatement{thm:main-classification}{Complete parameter classification}
  & --- & Not formalized\\
\BlueprintStatement{def:conditional-entropy-axioms}{Conditional-entropy axiom bundle}
  & --- & Not formalized\\
\BlueprintStatement{lem:embedding-lift}{Embedding lift of fixed-row monotonicity}
  & --- & Not formalized\\
\BlueprintStatement{cor:all-conditional-entropy-axioms}{Admissible candidates satisfy the complete axiom bundle}
  & --- & Not formalized\\

\end{longtable}
\normalsize

